\subsection*{Objetivo}
Integrar encoders, ponte H, sensor ultrassônico e comunicação sem fio (nRF24L01+) em um carrinho
que anda de forma autônoma, desvia de obstáculos e recebe comandos remotos.

\subsection{O que foi implementado, parcialmente implementado ou não implementado}

\paragraph{Implementado.}
\begin{itemize}[noitemsep]
    \item Reorganização do código em bibliotecas separadas (\texttt{encoder}, \texttt{motor\_control},
          \texttt{ultrasound}, \texttt{nrf24l01}), com parâmetros configuráveis via Kconfig
          (papel de transmissor ou receptor, distância mínima de 10\,cm, diâmetro da roda, bitola,
          pulsos por volta e tempo de parada).
    \item Controle de velocidade dos motores por PWM (canais do TPM2, período de 40\,\textmu s, ou seja, 25\,kHz),
          além dos comandos de frente, ré, esquerda, direita e parada.
    \item Sensor ultrassônico HC-SR04: pulso de \textit{trigger} de 10\,\textmu s gerado por PWM (período de 60\,ms),
          medição do \textit{echo} por interrupção com contador de ciclos, conversão para centímetros e
          notificação por semáforo quando a distância fica abaixo do limiar.
    \item Encoders com contadores atômicos e alvo de pulsos com notificação por semáforo,
          permitindo giros de ângulo arbitrário.
    \item Máquina de estados do carrinho (andar, desviar de obstáculo, parar), executada em \textit{threads}
          com \texttt{k\_poll}, semáforos e fila de mensagens. No desvio, o carrinho para, gira $90^\circ$ para a esquerda
          e verifica o caminho; se ainda houver obstáculo, gira $180^\circ$ para a direita e verifica novamente;
          caso persista, gira $90^\circ$ para a direita e para.
\end{itemize}

\paragraph{Parcialmente implementado.}
\begin{itemize}[noitemsep]
    \item \textbf{Enlace de rádio:} foram escritos o código de acesso por SPI ao nRF24L01+ (leitura e escrita de registradores,
          buffers e inicialização), o tratamento da interrupção IRQ, os papéis de transmissor e receptor
          escolhidos por Kconfig e o envio de comandos (\texttt{R}, \texttt{S}, \texttt{P}, \texttt{C}) digitados no terminal.
          No carrinho, apenas os comandos de \textit{run} e \textit{stop} têm efeito; \texttt{P} e \texttt{C}
          foram definidos, mas sem ação associada. Como a comunicação não funcionou, esta parte não pôde ser validada.
    \item \textbf{Calibração dos motores:} a função \texttt{motor\_calibrate} apenas fixa as velocidades em 80\% (esquerdo)
          e 100\% (direito), sem rotina automática. As constantes de iterações e tolerância de pulsos estão declaradas,
          mas não são utilizadas.
\end{itemize}

\paragraph{Não implementado / não funcional.}
\begin{itemize}[noitemsep]
    \item \textbf{A comunicação sem fio entre as placas não funcionou.} Consequentemente, o controle remoto
          do carrinho por rádio não pôde ser demonstrado.
\end{itemize}

\subsection{Dificuldades e soluções encontradas}

\paragraph{Principal dificuldade.}
O maior problema deste projeto foi o não funcionamento da comunicação sem fio entre as duas placas com
os módulos nRF24L01+. A falha não foi solucionada até a data da entrega.

\paragraph{Hipótese para a causa e possível solução.}
Como possível causa, suspeita-se de mau contato nos \textit{jumpers} usados para ligar os módulos de rádio
à placa, já que o nRF24L01+ é sensível a conexões instáveis nos sinais de SPI e de alimentação.
Essa suspeita é reforçada pelo fato de que, na Aula 2, a comunicação entre as placas funcionou
corretamente, com o controle remoto do LED pelo terminal. Esta hipótese ainda não foi confirmada por testes.

Uma possível solução seria montar o módulo de rádio em uma PCB, eliminando os \textit{jumpers} e deixando
as conexões de SPI, CE, IRQ e alimentação fixas. Essa alternativa foi apenas considerada e não será
implementada pelo grupo.
\paragraph{Outras dificuldades.}
\begin{itemize}
    \item \textbf{Sincronização entre interrupções e \textit{threads}} (encoders, ultrassom e controle do carrinho). \\
          \textbf{Solução:} contadores atômicos, semáforos de notificação, fila de mensagens e \texttt{k\_poll}
          para o carrinho reagir a obstáculos e a comandos sem espera ocupada.
    \item \textbf{Medição de distância com o ultrassom} (conversão de ciclos do relógio para microssegundos e depois para centímetros,
          e ruído em pulsos muito curtos). \\
          \textbf{Solução:} conversão com \texttt{k\_cyc\_to\_us\_floor32}, divisão por 58 e descarte de pulsos abaixo de
          uma duração mínima.
    \item \textbf{Motores com velocidades diferentes, fazendo o carrinho desviar da linha reta.} \\
          \textbf{Solução parcial:} controle por PWM com velocidades distintas por lado (80\% e 100\%),
          ajustadas manualmente, sem calibração automática.
\end{itemize}

\end{document}